\chapter{Experimental Validation}
\label{cap:experiments}

This chapter reports the experimental evaluation of the segmentation models, geometric extractors, and integrated navigation system. It covers offline evaluation using recorded images and, if completed, closed-loop field trials.

\section{Offline evaluation protocol}
\label{sec:offline_evaluation}

\subsection{Segmentation metrics}
\label{subsec:segmentation_metrics}

Segmentation performance is measured using intersection over union (IoU) for each semantic class. For a class $c$, IoU is defined as:
\begin{equation}
    \operatorname{IoU}_c=
    \frac{\operatorname{TP}_c}
         {\operatorname{TP}_c+\operatorname{FP}_c+\operatorname{FN}_c},
    \label{eq:iou}
\end{equation}
where true positives, false positives, and false negatives are computed against the manual reference masks.

Mean IoU (mIoU) averages the per-class IoU values. For the four-class configuration, mIoU includes background, plants, path, and covering material. Path IoU is the task-specific metric most relevant for navigation.

\subsubsection{Architecture comparison results}
\label{subsubsec:architecture_results}

The architecture comparison results reported in Section~\ref{subsec:c4_architecture_results} demonstrate that LR-ASPP achieves path IoU of 0.9546 with inference time of 10.41~ms, compared to FCN (0.9621, 58.24~ms) and DeepLabV3 (0.9651, 86.87~ms). The 1.2 percentage point mIoU deficit is offset by a 5--8$\times$ reduction in inference latency.

\subsubsection{Supervision strategy results}
\label{subsubsec:supervision_results}

% TODO: Complete after provenance verification
% - Compare manual 4-class vs SAM 4-class
% - Compare manual 4-class vs SAM 3-class (plants+path)
% - Compare manual vs SAM binary (path only)
% - Report on full test set and operational subset

\subsection{Geometric reference accuracy}
\label{subsec:geometric_accuracy}

Geometric reference quality is assessed by comparing Method~3 and Method~4 outputs on the same evaluation masks. Metrics include:
\begin{itemize}
\item Reference line parameters (slope and intercept)
\item Lateral deviation from manual reference (if available)
\item Rejection rate (percentage of images where extraction fails validation)
\item Computational time
\end{itemize}

% TODO: Add geometric comparison results
% - Method 3 vs Method 4 on lettuce test set
% - Effect of ROI, strip count, and parameter choices
% - Success rate and failure modes

\section{Field validation}
\label{sec:field_validation}

\subsection{Test conditions}
\label{subsec:test_conditions}

% TODO: If field trials completed, describe:
% - Test location and crop type
% - Row length and width
% - Lighting and weather conditions
% - Number of trials

\subsection{Closed-loop navigation results}
\label{subsec:navigation_results}

% TODO: If field trials completed, report:
% - Lateral tracking error
% - Success rate (percentage of row length navigated)
% - Failure modes (loss of reference, collisions, controller instability)
% - Comparison with baseline methods (if applicable)

\subsection{Qualitative observations}
\label{subsec:qualitative}

% TODO: If field trials completed, discuss:
% - Robustness to gaps and occlusions
% - Behavior under shadow and lighting changes
% - Limitations and edge cases

\section{Discussion}
\label{sec:discussion}

\subsection{Segmentation performance}
\label{subsec:segmentation_discussion}

The architecture comparison demonstrates that LR-ASPP provides a practical accuracy--speed trade-off for robotic deployment. Its path IoU of 0.9546 on manually annotated lettuce is sufficient for geometric extraction, while its inference time supports real-time operation.

The supervision strategy comparison addresses the effort required for dataset preparation. SAM~3-assisted annotation reduces manual labeling time, but its quality and the need for inspection remain to be quantified systematically.

% TODO: Expand discussion with supervision results

\subsection{Geometric extraction performance}
\label{subsec:geometric_discussion}

Method~3 and Method~4 differ in their association strategy and validation checks. Their relative performance depends on mask quality, row structure, and parameter tuning.

% TODO: Expand with geometric comparison results

\subsection{Limitations}
\label{subsec:limitations}

The experiments are subject to several limitations:
\begin{itemize}
\item Segmentation evaluation uses manually annotated lettuce as ground truth. Performance on other crops is inferred from multi-crop training but not directly validated against manual references for tomato or cucumber.
\item Geometric extractors assume locally straight boundaries. They do not handle curves or complex row patterns.
\item Offline evaluation does not establish closed-loop navigation performance. Segmentation IoU and geometric reference quality are necessary but not sufficient for successful field operation.
\item If field trials were not completed, the integrated system's behavior under real navigation conditions remains unvalidated.
\end{itemize}

% TODO: Update limitations based on completed experiments
